\documentclass[conference]{IEEEtran}
\usepackage{amsmath,amsfonts}
\usepackage{array}
\usepackage{booktabs}
\usepackage[caption=false,font=footnotesize]{subfig}
\usepackage{graphicx}
\usepackage{cite}
\usepackage{url}

\hyphenation{op-tical net-works semi-conduc-tor}

\begin{document}

\title{FX-RV32-X: A Multi-Bank Shadow Register Architecture\\ for Priority-Based Nested Interrupts with \\ Constant 2-Cycle Latency}

\author{\IEEEauthorblockN{Fengxin Yi}
\IEEEauthorblockA{Hangzhou International Innovation Institute\\
BeiHang University\\
Hangzhou, China\\
Email: 1596215367@buaa.edu.cn}}

\maketitle

\begin{abstract}
Hard real-time embedded systems require both low interrupt latency and priority-based preemptive nesting, yet existing RISC-V solutions either lack nesting support or incur serial context-save overhead proportional to nesting depth.
We present FX-RV32-X, an extension to the FX-RV32 processor that replaces the single shadow register bank with an $N$-bank array managed by a purely combinational Bank Controller.
The Bank Controller performs priority preemption judgment, automatic Bank allocation and release, overflow handling, and tail-chaining optimization—all within the 2-cycle interrupt entry window, adding zero delay to the critical path.
Interrupt latency remains fixed at 2 clock cycles regardless of nesting depth.
We implemented FX-RV32-X in 55 nm CMOS (four configurations, 1--8 Banks) and on a Xilinx Kintex-7 FPGA.
The 4-Bank ASIC configuration occupies 62.4 kGE with a per-Bank marginal cost of $\sim$9.5 kGE and zero timing penalty (critical path variation $<$0.2\%).
FPGA implementation achieves 160--178 MHz core-logic Fmax across all configurations.
Direct measurement of the OpenE902 CLIC core confirms FX-RV32-X completes three-level nested interrupt entry and exit in 15 cycles—1.5$\times$ lower than E902 (22 cycles) and 4.8$\times$ lower than CLIC hardware stacking, while saving 31 registers per level versus zero in E902.
\end{abstract}

\section{Introduction}

Hard real-time systems in automotive, industrial control, and aerospace domains must handle multiple asynchronous interrupt sources with different urgency levels.
In servo motor control, the current loop (50 $\mu$s period) must preempt the velocity loop (200 $\mu$s); in electric power steering, torque-overload protection must preempt angle-sampling interrupts.
These scenarios demand two capabilities: priority-based preemptive nesting, and interrupt latency that is both low and constant regardless of nesting depth.

Our prior work, FX-RV32~\cite{ref1}, achieved a fixed 2-cycle interrupt latency using a single set of 31 hardware shadow registers for parallel context save.
However, the single-bank design cannot support interrupt nesting—a higher-priority interrupt overwrites the previous shadow register content, causing irreversible context loss.

Existing solutions exhibit a fundamental tradeoff.
The CLIC hardware stacking scheme~\cite{ref2} provides independent per-level storage but its serial push/pop operation scales linearly with the number of saved registers per level (11 cycles for 10 registers).
The Sophon processor~\cite{ref3} offers single-cycle parallel save via a custom \texttt{snapreg} instruction but explicitly does not support nesting.
The OpenE902 core~\cite{ref4} implements CLIC vectored interrupts but provides no hardware context save at all—software must execute \texttt{sw} instructions in the ISR prologue.
ARM NVIC~\cite{ref5} performs automatic hardware stacking but is proprietary, saves only 8 registers, and requires 12 cycles per entry.

This paper presents FX-RV32-X, which resolves this tradeoff by extending the single shadow register bank to $N$ independent banks managed by a hardware Bank Controller.
The key innovations are: (1) a multi-bank shadow register array enabling $N-1$ levels of interrupt nesting; (2) a purely combinational Bank Controller that makes preemption, allocation, overflow, and tail-chain decisions in zero additional cycles; (3) configurable overflow policies for resource-exhaustion scenarios; and (4) complete software transparency requiring no ISA extensions.

\section{Multi-Bank Shadow Register Architecture}

\subsection{Overall Architecture}

FX-RV32-X extends the baseline five-stage FX-RV32 pipeline by replacing the single shadow register bank with an $N \times 31$ multi-bank array integrated within the register file, and adding a Bank Controller module.
Fig.~\ref{fig_arch} shows the architecture.

\begin{figure}[!t]
\centering
\includegraphics[width=\columnwidth]{fig1_v2.png}
\caption{Overall architecture. The Bank Controller (highlighted) is purely combinational. The multi-bank shadow array resides within the register file.}
\label{fig_arch}
\end{figure}

The Bank Controller receives the new interrupt's priority from the interrupt controller, the currently-serviced interrupt's priority from the interrupt pipeline, and the current Bank pointer value.
It outputs four control signals via combinational logic: \texttt{allow\_nesting} (enables preemption), \texttt{bank\_full} (blocks preemption when all Banks occupied), \texttt{tail\_chain\_detect} (suppresses context restore on MRET+pending), and \texttt{degradation\_reuse} (enables Bank overwrite in overflow Policy~B).

\subsection{Bank Controller: Preemption and Allocation}

The preemption judgment is a single combinational expression:

\begin{equation}
\text{allow\_nesting} = (\text{new\_prio} > \text{cur\_prio}) \land (\text{bank\_ptr} < N) \land \lnot\text{degrad\_flag}
\label{eq:preempt}
\end{equation}

When the interrupt pipeline accepts an interrupt, all register updates occur simultaneously at the same clock edge: \texttt{interrupt\_taken} $\leftarrow$ 1 (1-cycle pulse), \texttt{shadow\_save} $\leftarrow$ 1, \texttt{interrupt\_flush} $\leftarrow$ 1 (clears ID/EX, EX/MEM, MEM/WB but preserves IF/ID), CSR registers (mepc, mcause, mstatus) are written, and \texttt{bank\_ptr} increments if nesting is allowed.

The register file uses a three-level write priority (Eq.~\ref{eq:priority}) preserved from the baseline design:

\begin{equation}
\text{shadow\_restore} > \text{normal WB} > \text{shadow\_save}
\label{eq:priority}
\end{equation}

This ensures that any instruction completing WB before the interrupt entry finishes its write before the shadow snapshot is taken, and that shadow\_restore on MRET completely overwrites ISR-internal register modifications.

\subsection{Constant 2-Cycle Latency}

Fig.~\ref{fig_timing} illustrates the interrupt entry timing.
At T1, the interrupt controller asserts \texttt{intr\_pending} combinatorially; \texttt{intr\_take\_now} drives \texttt{next\_pc} to the handler address in the same cycle.
At T2$\uparrow$, all acceptance operations complete simultaneously.
By T4$\uparrow$, the first ISR instruction enters the EX stage—exactly 2 cycles from acceptance.
This timing is identical for the first interrupt and every subsequent nesting level.

\begin{figure}[!t]
\centering
\includegraphics[width=\columnwidth]{fig_timing_v2.png}
\caption{Constant 2-cycle interrupt entry timing. The unconditional acceptance mechanism (intr\_take\_now) enables PC redirection in the same cycle as the request. All register operations occur at the same clock edge (T2$\uparrow$).}
\label{fig_timing}
\end{figure}

\subsection{Tail-Chaining Optimization}

In conventional handling, an MRET followed immediately by a new interrupt incurs a redundant ``restore-then-save'' sequence: the MRET triggers a context restore from the current Bank, then the new interrupt triggers a context save back to the same Bank.
When the Bank Controller detects \texttt{id\_ex\_mret} $\land$ \texttt{intr\_pending} combinatorially, it suppresses the shadow\_restore pulse and holds \texttt{bank\_ptr} unchanged.
The new interrupt directly reuses the current Bank, compressing the MRET-to-next-ISR overhead from 3 cycles to 2 cycles.

\subsection{Overflow Handling}

When all $N$ Banks are occupied, the Bank Controller provides two configurable policies via the \texttt{OVERFLOW\_POLICY} synthesis parameter:

\begin{itemize}
\item \textbf{Hard Limit} (POL=0, default): Asserts \texttt{bank\_full}, blocking the new interrupt and preserving all preempted contexts. The interrupt remains pending until an MRET frees a Bank.
\item \textbf{Degradation Reuse} (POL=1): Overwrites Bank[$N-1$] (the deepest nesting level), guaranteeing the highest-priority interrupt is always admitted within 2 cycles. The overwritten ISR's MRET chain is broken; the remaining $N-1$ levels are unaffected.
\end{itemize}

\section{Experimental Results}

\subsection{Setup}

FX-RV32-X was synthesized with Synopsys Design Compiler (SMIC 55 nm, compile\_ultra) for $N = 1, 2, 4, 8$ Banks at 200 MHz.
FPGA implementation used Vivado 2022.2 on a Xilinx Kintex-7 XC7K325T (-2 speed grade).
Functional validation comprised ten test programs (single interrupt, two-level nesting, three-level nesting, context integrity, both overflow policies, and tail-chaining) verified in ModelSim simulation.
For competitive comparison, we synthesized and simulated the OpenE902 CLIC core~\cite{ref4} in iverilog to directly measure its interrupt and nested-interrupt latency.

\subsection{ASIC Area Scaling}

Table~\ref{tab:area} and Fig.~\ref{fig_area} present the synthesis results.
The Banks=1 configuration adds 3.63 kGE of multi-bank infrastructure (Bank Controller $\sim$0.02 kGE, MUX tree and bank\_ptr logic $\sim$3.61 kGE) over the original single-bank FX-RV32 (32.37 kGE).
Beyond the first Bank, area scales at $\sim$9.5 kGE per additional Bank, comprising 31$\times$32-bit flip-flops ($\sim$7.0 kGE) and MUX tree expansion ($\sim$2.5 kGE).
The critical path remains at 4.95--4.96 ns (variation $<$ 0.2\%), confirming the Bank Controller does not lie on any timing-critical path.

\begin{table}[!t]
\caption{ASIC Synthesis Results (55 nm, 200 MHz)}
\label{tab:area}
\centering
\footnotesize
\setlength{\tabcolsep}{3pt}
\begin{tabular}{@{}lcccc@{}}
\toprule
\textbf{$N$} & \textbf{Area (kGE)} & \textbf{$\Delta$ (kGE)} & \textbf{Power (mW)} & \textbf{T$_{\text{crit}}$ (ns)} \\
\midrule
1 & 36.00 & --     & 6.07 & 4.95 \\
2 & 42.99 & +6.99  & 7.70 & 4.96 \\
4 & 62.41 & +19.42 & 11.54 & 4.96 \\
8 & 100.20 & +37.79 & 19.38 & 4.96 \\
\bottomrule
\end{tabular}
\end{table}

\begin{figure}[!t]
\centering
\includegraphics[width=\columnwidth]{fig_area_v2.png}
\caption{ASIC area scaling with Bank count.}
\label{fig_area}
\end{figure}

\subsection{FPGA Implementation}

Table~\ref{tab:fpga} summarizes the post-route FPGA results.
FF count grows linearly at 986--1,005 per additional Bank, matching the theoretical 31$\times$32 = 992 FFs.
LUT growth is sub-linear and saturates after $N = 4$ (6,829 $\rightarrow$ 6,838 LUTs for $N = 4 \rightarrow 8$), as FPGA LUT6/F7MUX/F8MUX primitives absorb the additional MUX tree depth.
Setup WNS shows no monotonic degradation across configurations ($-0.612$ to $-1.235$ ns, $\pm$0.3 ns variation from P\&R noise).
Core-logic Fmax ranges from 160--178 MHz, limited by the base CPU design, not the Bank mechanism.

\begin{table}[!t]
\caption{FPGA Implementation Results (Kintex-7, Post-Route)}
\label{tab:fpga}
\centering
\footnotesize
\setlength{\tabcolsep}{3pt}
\begin{tabular}{@{}lcccc@{}}
\toprule
\textbf{$N$} & \textbf{LUTs} & \textbf{FFs} & \textbf{WNS (ns)} & \textbf{Fmax (MHz)} \\
\midrule
1 & 4,356 & 3,917  & -0.748 & 174.0 \\
2 & 4,855 & 4,905  & -1.125 & 163.3 \\
4 & 6,829 & 6,914  & -0.612 & 178.2 \\
8 & 6,838 & 10,859 & -1.235 & 160.4 \\
\bottomrule
\end{tabular}
\end{table}

\subsection{Nested Interrupt Latency}

We measured entry and exit latency using a three-level nested interrupt test program (Software $\rightarrow$ Timer $\rightarrow$ GPIO) at 200 MHz in simulation.
Table~\ref{tab:nested} reports the results.

\begin{table}[!t]
\caption{Measured Nested Interrupt Latency (3-Level Test)}
\label{tab:nested}
\centering
\footnotesize
\setlength{\tabcolsep}{3pt}
\begin{tabular}{@{}lccc@{}}
\toprule
\textbf{Level} & \textbf{Entry} & \textbf{Exit} & \textbf{Cumulative} \\
\midrule
1 (SW, ID=3)    & 2 cycles & 3 cycles & 5 cycles \\
2 (Timer, ID=7) & 2 cycles & 3 cycles & 10 cycles \\
3 (GPIO, ID=11) & 2 cycles & 3 cycles & 15 cycles \\
\bottomrule
\end{tabular}
\end{table}

Entry latency is exactly 2 cycles at every nesting level.
The cumulative round-trip overhead of 15 cycles (75 ns at 200 MHz) is independent of ISR body length.

\subsection{Comparison with CLIC-Based Implementations}

Table~\ref{tab:compare} compares FX-RV32-X with two CLIC-based approaches: the OpenE902 core~\cite{ref4} measured in this work, and the CLIC hardware stacking scheme~\cite{ref2} (literature data, architectural extrapolation for nesting).

\begin{table}[!t]
\caption{Nested Interrupt Overhead Comparison}
\label{tab:compare}
\centering
\footnotesize
\setlength{\tabcolsep}{3pt}
\begin{tabular}{@{}lccc@{}}
\toprule
\textbf{Metric} & \textbf{FX-RV32-X} & \textbf{E902 CLIC}$^a$ & \textbf{CLIC Stack~\cite{ref2}} \\
\midrule
Entry latency (1 level)     & 2 cycles   & 3 cycles$^a$  & 13 cycles \\
Entry latency (nested)      & 2 cycles   & 5 cycles$^a$  & 13 cycles$^b$ \\
Exit latency                & 3 cycles   & 3 cycles      & $\sim$11 cycles$^b$ \\
3-level cumulative          & 15 cycles  & 22 cycles     & 72 cycles$^b$ \\
Hardware context save       & 31 regs, 0 cyc. & None (software) & 2--10 regs, 11 cyc. \\
\bottomrule
\multicolumn{4}{@{}p{0.95\columnwidth}}{\footnotesize $^a$Measured in iverilog simulation (this work). E902 has zero hardware register save; context must be saved by software \texttt{sw} instructions ($\sim$32 cycles/level for RV32E). $^b$Estimated from per-level serial push/pop cost~\cite{ref2}; ref.~\cite{ref2} reports only single-interrupt latency.} \\
\end{tabular}
\end{table}

FX-RV32-X achieves 1.5$\times$ lower cumulative overhead than E902 (15 vs.\ 22 cycles) while saving 31 registers per level versus zero hardware-saved registers in E902.
When E902's mandatory software context save ($\sim$32 cycles/level) is included, the advantage exceeds 7$\times$ per nesting level.
Against the CLIC hardware stacking scheme, the reduction is 4.8$\times$.

A key finding from the E902 analysis is that the core \emph{does not implement hardware register stacking}—its RTL contains no push, pop, or shadow register logic.
This confirms that the per-level $+$2 cycle nesting penalty observed in E902 (3$\rightarrow$5 cycles entry) arises purely from CLIC re-arbitration, which FX-RV32-X's purely combinational Bank Controller eliminates.

\subsection{Functional Validation}

All ten functional validation tests passed, including three-level nesting with full 31-register context integrity verification (all x1--x31 correctly restored to pre-interrupt values after nested preemption return).
The tail-chaining test confirmed a 1-cycle saving per tail-chain event (29 cycles from GPIO to Timer entry vs.\ 30 cycles under the normal path).

\section{Conclusion}

FX-RV32-X extends the FX-RV32 processor with a multi-bank shadow register array and a purely combinational Bank Controller, enabling priority-based nested interrupts while preserving the baseline 2-cycle constant interrupt latency at every nesting level.
The architecture achieves zero-cycle context-save overhead (parallel lock within the entry window), zero-cycle preemption decision latency (combinational priority comparison), and zero-measurable timing impact on the critical path across 1--8 Bank configurations in both 55 nm ASIC and 28 nm FPGA implementations.
Direct competitive measurement of OpenE902 confirms FX-RV32-X provides the only RISC-V solution that simultaneously supports multi-level hardware nesting, complete 31-register parallel context save, and constant 2-cycle latency.

\begin{thebibliography}{5}

\bibitem{ref1}
F.~Yi, ``FX-RV32: A Lightweight, Deterministic and Low Latency RISC-V Processor for Hard Real-Time Embedded Systems,'' \textit{IEEE Trans. Very Large Scale Integr. (VLSI) Syst.}, submitted for publication.

\bibitem{ref2}
B.~Mao, N.~Tan, T.~Chong, and L.~Li, ``A CLIC extension based fast interrupt system for embedded RISC-V processors,'' in \textit{Proc. 6th Int. Conf. Integr. Circuits Microsystems (ICICM)}, Nanjing, China, 2021, pp.~109--113.

\bibitem{ref3}
Z.~Huang, X.~Chen, F.~Gao, R.~Li, X.~Wu, and F.~Zhang, ``Sophon: A time-repeatable and low-latency architecture for embedded real-time systems based on RISC-V,'' \textit{IEEE Trans. Very Large Scale Integr. (VLSI) Syst.}, vol.~33, no.~1, pp.~221--233, Jan.~2025.

\bibitem{ref4}
T-Head Semiconductor, ``OpenE902 -- an open-source RISC-V core,'' 2020. [Online]. Available: \url{https://github.com/T-head-Semi/opene902}

\bibitem{ref5}
J.~Yiu, \textit{The Definitive Guide to ARM Cortex-M3 and Cortex-M4 Processors}, 3rd~ed. Newnes, 2013.

\end{thebibliography}

\end{document}
